% To compile using latexmk on the command line, run the following: 
% latexmk -pdf methods.tex

\documentclass[11pt]{article}

% ---------- Page formatting ----------
\usepackage{geometry}
\geometry{margin=1in}

\usepackage{lmodern}
\usepackage[T1]{fontenc}
\usepackage{setspace}
\onehalfspacing

% ---------- Graphics ----------
\usepackage{graphicx}
\usepackage{float}

% ---------- Math ----------
\usepackage{amsmath}
\usepackage{amssymb}
\usepackage{bm}

% ---------- ACS citations ----------
% \usepackage{achemso}
% \setkeys{acs}{articletitle=true}

% ---------- Title ----------
\title{\textbf{Computational Identification of Unknown Organic Contaminants and Molecules in Remote Environments via Molecular Networking}}
\author{Raelyn Brooks}
\date{\today}

\begin{document}

\maketitle

\nocite{*}

\begin{doublespace}

\section{Goal}
    This project aims to characterize surface water samples from remote high-altitude environments by identifying both known and unknown organic contaminants using high-resolution mass spectrometry (HRMS).
    While database-driven approaches with lists of known contaminants enable identification of some contaminants, I hypothesize that a significant fraction of uncharacterized molecules can be identified through non-targeted analysis, particularly molecular networking.
    Together, these approaches will reveal the identity and patterns of contaminants in these remote regions.

\section{Methods}
    \subsection{Sample Collection and Preparation}
        Surface water samples collected from remote high-altitude environments—including Mt. Everest, Mt. Kilimanjaro, the Tibetan Plateau, and Qinghai Lake—will be analyzed for organic contaminants. 
        These samples represent integrated records of atmospheric deposition.
        Prior to analysis, samples will be filtered to remove particulate matter and stored under controlled conditions to minimize degradation or contamination. 
        Solid-phase extraction (SPE) will be used to concentrate dissolved organic compounds and improve detection sensitivity for trace-level contaminants. 
        Extracted samples will then be reconstituted in an appropriate solvent compatible with liquid chromatography–mass spectrometry (LC-HRMS) analysis.

    \subsection{High-Resolution Mass Spectrometry Data Acquisition}
        All samples will be analyzed using liquid chromatography coupled with high-resolution mass spectrometry (LC-HRMS). 
        An iterative data-dependent acquisition (DDA) approach will be employed, in which each sample is injected and analyzed multiple times.
        In each successive run, previously fragmented precursor ions will be excluded using dynamic exclusion lists, allowing for increased fragmentation coverage of lower-abundance compounds. 
        This iterative DDA strategy improves the detection of unique spectral features and enhances MS/MS coverage across complex environmental samples.
        Mass spectra will be collected in both MS1 (full scan) and MS/MS modes to obtain accurate mass measurements and fragmentation patterns necessary for compound identification.

    \subsection{Data Processing and Feature Detection}
    Raw LC-HRMS data will be processed using software such as MS-DIAL. 
    Data processing steps will include peak detection, deconvolution, alignment across samples, and normalization.
    Detected features will be characterized by their mass-to-charge ratio (m/z), retention time, and MS/MS fragmentation spectra. 
    Feature tables will be generated to represent the full set of detected compounds across all samples.
    Quality control steps, including blank subtraction and filtering of low-confidence features, will be applied to remove background contaminants and noise.

    \subsection{Suspect Screening for Known Contaminants}
        To identify known environmental contaminants, suspect screening will be performed using curated chemical databases, such as the Agilent environmental contaminant database.
        Detected features will be matched against database entries using accurate mass measurements and MS/MS spectral similarity. 
        Candidate identifications will be assigned based on established confidence levels (e.g., Schymanski levels), incorporating mass accuracy, isotopic patterns, and fragmentation matching.
        This approach will enable the annotation of known compounds such as pharmaceuticals, pesticides, and industrial chemicals, providing a baseline characterization of anthropogenic contamination in the samples.

    \subsection{Molecular Networking and Non-Targeted Analysis}
        To characterize unknown compounds not identified through suspect screening, molecular networking will be performed using MS/MS spectral similarity.
        Spectral data will be uploaded to platforms such as GNPS (Global Natural Products Social Molecular Networking), where pairwise spectral comparisons will be used to construct molecular networks. 
        In these networks, nodes represent individual compounds, and edges represent spectral similarity relationships.
        Clusters of related compounds will be analyzed to identify structural relationships between known and unknown molecules. 
        Unknown compounds that cluster with annotated compounds can be tentatively characterized based on shared fragmentation patterns and mass differences.
        Additional annotation tools (e.g., MolNetEnhancer) may be used to integrate chemical class information and improve structural interpretation.

    \subsection{Structural Interpretation and Data Analysis}
        Mass differences between related compounds in molecular networks will be examined to infer potential chemical transformations (e.g., oxidation, methylation, or degradation products). 
        Fragmentation patterns will be interpreted to propose tentative structures for unknown compounds.
        Comparative analysis across sampling locations will be conducted to identify trends in contaminant distribution and to assess the influence of long-range atmospheric transport. 
        Differences in compound profiles may reveal region-specific contamination patterns or shared global sources.

\end{doublespace}

% ---------- References ----------
\newpage
\bibliographystyle{plain}
\bibliography{bibliography.bib}

\end{document}